\documentclass[11pt]{article}
\usepackage[margin=1in]{geometry}
\usepackage{hyperref}
\usepackage{listings}
\usepackage{xcolor}
\usepackage{longtable}
\usepackage{booktabs}
\usepackage{enumitem}
\usepackage{amsmath}

\lstset{
    basicstyle=\ttfamily\small,
    backgroundcolor=\color{gray!10},
    frame=single,
    framerule=0.5pt,
    breaklines=true,
    showstringspaces=false,
    columns=flexible
}

\title{FEMM Heat Flow File Formats\\[0.5em]
\large Input (\texttt{.feh}) and Output (\texttt{.anh}) Specifications}
\author{D.\ Meeker}
\date{\today}

\begin{document}
\maketitle
\tableofcontents
\newpage

% ======================================================================
\section{Overview}

The FEMM heat flow solver (\texttt{hsolv}) reads a problem definition
from a \texttt{.feh} file and writes the solution to a \texttt{.anh}
file.  The \texttt{.anh} file begins with a verbatim copy of the input
\texttt{.feh} file, followed by a \texttt{[Solution]} section
containing the finite element mesh and computed temperature values.

All floating-point values are written with 17 significant digits
(\verb|%.17g| format).  Fields within a line are separated by tabs or
spaces.  Section headers use square brackets (\verb|[Tag]|); property
block delimiters use angle brackets (\verb|<BeginBlock>| /
\verb|<EndBlock>|).  String comparisons during parsing are
case-insensitive.

% ======================================================================
\section{The \texttt{.feh} Input File}

% ----------------------------------------------------------------------
\subsection{Global Problem Parameters}

Each parameter occupies one line in the format
\verb|[Tag] = value|.  All are optional; defaults are applied if
missing.

\begin{longtable}{lllp{5.5cm}}
\toprule
Tag & Type & Default & Description \\
\midrule
\endhead
\verb|[Format]|      & int     & 1          & File format version \\
\verb|[Frequency]|   & int     & 0          & Always 0 (legacy field) \\
\verb|[Precision]|   & double  & 1e-8       & Nonlinear convergence tolerance \\
\verb|[MinAngle]|    & double  & 30         & Minimum mesh element angle (degrees) \\
\verb|[DoSmartMesh]| & int     & 1          & Adaptive meshing (0=off, 1=on) \\
\verb|[Depth]|       & double  & 1.0        & Extrusion depth (planar problems) \\
\verb|[LengthUnits]| & string  & millimeters & inches, millimeters, centimeters, meters, mils, microns \\
\verb|[ProblemType]| & string  & planar     & planar or axisymmetric \\
\verb|[Coordinates]| & string  & cartesian  & cartesian or polar \\
\verb|[Comment]|     & string  & (empty)    & Problem description (quoted; \verb|\n| = newline) \\
\verb|[ACSolver]|    & int     & 0          & Always 0 (legacy field) \\
\verb|[dT]|          & double  & 0          & Time step for transient analysis (0 = steady-state) \\
\verb|[PrevSoln]|    & string  & (empty)    & Path to previous solution for transient initialization (quoted) \\
\bottomrule
\end{longtable}

The \verb|[dT]| and \verb|[PrevSoln]| fields enable transient thermal
analysis.  When \verb|dT| $> 0$, the solver performs a single time step
of duration \verb|dT|, using the temperature field from
\verb|PrevSoln| as the initial condition.  Steady-state analysis
corresponds to \verb|dT| $= 0$.

\subsubsection{Axisymmetric External Region}

Written only when \verb|ProblemType| is axisymmetric and both
\verb|extRo| and \verb|extRi| are nonzero:

\begin{longtable}{llp{8cm}}
\toprule
Tag & Type & Description \\
\midrule
\verb|[extZo]| & double & Outer Z coordinate \\
\verb|[extRo]| & double & Outer radius \\
\verb|[extRi]| & double & Inner radius \\
\bottomrule
\end{longtable}

% ----------------------------------------------------------------------
\subsection{Point (Nodal) Properties}

\begin{lstlisting}
[PointProps] = <count>
<BeginPoint>
  <PointName> = "name"
  <Tp> = 0        prescribed temperature (K)
  <qp> = 0        point heat source (W)
<EndPoint>
... repeated count times
\end{lstlisting}

% ----------------------------------------------------------------------
\subsection{Boundary Properties}

\begin{lstlisting}
[BdryProps] = <count>
<BeginBdry>
  <BdryName> = "name"
  <BdryType> = 0
  <Tset> = 0          fixed temperature (K)
  <qs> = 0            prescribed heat flux (W/m^2)
  <Tinf> = 0          ambient temperature for convection (K)
  <h> = 0             convection coefficient (W/m^2-K)
  <beta> = 0          radiation coefficient (W/m^2-K^4)
  <TinfRad> = 0       radiation reference temperature (K)
<EndBdry>
\end{lstlisting}

Boundary types:
\begin{itemize}[nosep]
  \item 0 = fixed temperature (Dirichlet): $T = T_\text{set}$
  \item 1 = prescribed heat flux (Neumann): $\hat{n} \cdot F = q_s$
  \item 2 = convection: $\hat{n} \cdot F = h(T - T_\infty) + q_s$
  \item 3 = radiation + convection:
        $\hat{n} \cdot F = h(T - T_\infty) + \beta(T^4 - T_{\infty,\text{rad}}^4) + q_s$
  \item 4 = periodic
  \item 5 = antiperiodic
\end{itemize}

If \verb|TinfRad| is not present in the file, it defaults to the value
of \verb|Tinf| (backward compatibility).

% ----------------------------------------------------------------------
\subsection{Material (Block) Properties}

\begin{lstlisting}
[BlockProps] = <count>
<BeginBlock>
  <BlockName> = "name"
  <Kx> = 1          thermal conductivity, x-direction (W/m-K)
  <Ky> = 1          thermal conductivity, y-direction (W/m-K)
  <Kt> = 0          volumetric heat capacity (J/m^3-K)
  <qv> = 0          volumetric heat generation (W/m^3)
  <TKPoints> = N
    T_0  K_0         temperature-conductivity data (K, W/m-K)
    T_1  K_1         ... repeated N times
    ...
<EndBlock>
\end{lstlisting}

When \verb|TKPoints| $> 0$, the thermal conductivity is nonlinear:
$K(T)$ is piecewise-linear interpolated from the T-K curve.  Below the
first data point, the first conductivity value is used; above the last
data point, the last value is used.  The solver iterates until
convergence when nonlinear materials or radiation BCs are present.

% ----------------------------------------------------------------------
\subsection{Conductor Properties}

Conductors enforce equipotential (isothermal) constraints or prescribe
total heat flow on surfaces, analogous to circuits in the magnetics
solver.

\begin{lstlisting}
[ConductorProps] = <count>
<BeginConductor>
  <ConductorName> = "name"
  <Tc> = 0           prescribed temperature (K)
  <qc> = 0           prescribed total heat flow (W)
  <ConductorType> = 0
<EndConductor>
\end{lstlisting}

Conductor types:
\begin{itemize}[nosep]
  \item 0 = prescribed heat flow ($q = q_c$; solver computes $T$)
  \item 1 = prescribed temperature ($T = T_c$; solver computes $q$)
\end{itemize}

% ----------------------------------------------------------------------
\subsection{Geometry}

\subsubsection{Nodes}

\begin{lstlisting}
[NumPoints] = <count>
x    y    propIdx    inGroup    conductorIdx    [lfs]
\end{lstlisting}

\begin{itemize}[nosep]
  \item \texttt{x}, \texttt{y}: coordinates (double)
  \item \texttt{propIdx}: index into point property list (1-based; 0 = none)
  \item \texttt{inGroup}: group number (int)
  \item \texttt{conductorIdx}: index into conductor list (1-based; 0 = none)
  \item \texttt{lfs}: local feature size constraint (double, optional; $-1$ = auto)
\end{itemize}

\subsubsection{Line Segments}

\begin{lstlisting}
[NumSegments] = <count>
n0    n1    maxSideLength    propIdx    hidden    inGroup    conductorIdx
\end{lstlisting}

\begin{itemize}[nosep]
  \item \texttt{n0}, \texttt{n1}: endpoint node indices (0-based)
  \item \texttt{maxSideLength}: mesh size constraint (double; $-1$ = auto)
  \item \texttt{propIdx}: boundary property index (1-based; 0 = none)
  \item \texttt{hidden}: 0 = visible, 1 = suppressed from mesh
  \item \texttt{inGroup}: group number
  \item \texttt{conductorIdx}: conductor index (1-based; 0 = none)
\end{itemize}

\subsubsection{Arc Segments}

\begin{lstlisting}
[NumArcSegments] = <count>
n0  n1  arcAngle  maxSideLength  propIdx  hidden  inGroup  conductorIdx  mySideLength
\end{lstlisting}

\begin{itemize}[nosep]
  \item \texttt{n0}, \texttt{n1}: endpoint node indices (0-based)
  \item \texttt{arcAngle}: arc span (degrees)
  \item \texttt{maxSideLength}: mesh constraint (double)
  \item \texttt{propIdx}: boundary property index (1-based; 0 = none)
  \item \texttt{hidden}: 0 = visible, 1 = suppressed
  \item \texttt{inGroup}: group number
  \item \texttt{conductorIdx}: conductor index (1-based; 0 = none)
  \item \texttt{mySideLength}: computed side length (0 = auto)
\end{itemize}

\subsubsection{Holes}

\begin{lstlisting}
[NumHoles] = <count>
x    y    inGroup
\end{lstlisting}

\subsubsection{Block Labels}

\begin{lstlisting}
[NumBlockLabels] = <count>
x    y    blockIdx    meshSize    inGroup    extFlags
\end{lstlisting}

\begin{itemize}[nosep]
  \item \texttt{x}, \texttt{y}: label position (double)
  \item \texttt{blockIdx}: material index (1-based; 0 = unassigned)
  \item \texttt{meshSize}: element size (double; $\leq 0$ = auto).
        Stored as diameter; converted to area: $A = \pi d^2/4$
  \item \texttt{inGroup}: group number
  \item \texttt{extFlags}: bit 0 = external region, bit 1 = default region
\end{itemize}

% ----------------------------------------------------------------------
\subsection{Differences from Magnetics and Electrostatics Formats}

\begin{itemize}
  \item \texttt{[dT]} and \texttt{[PrevSoln]} enable transient analysis
        (unique to heat flow).
  \item Material properties use thermal conductivity ($K_x$, $K_y$),
        volumetric heat capacity ($K_t$), and heat generation ($q_v$).
  \item Nonlinear materials use T-K curves (temperature vs.\
        conductivity) instead of B-H curves.
  \item Boundary conditions include convection ($h$, $T_\infty$) and
        radiation ($\beta$, $T_{\infty,\text{rad}}$), which are
        physics-specific to heat transfer.
  \item Point properties use temperature $T_p$ and point heat
        source $q_p$.
\end{itemize}


% ======================================================================
\section{The \texttt{.anh} Output File}

The \texttt{.anh} file begins with a verbatim copy of the input
\texttt{.feh} file (all sections described in Section~2), followed by
the solver's solution data.

% ----------------------------------------------------------------------
\subsection{Solution Header}

Immediately after the echoed \texttt{.feh} content:

\begin{lstlisting}
[Solution]
\end{lstlisting}

% ----------------------------------------------------------------------
\subsection{Mesh Nodes}

\begin{lstlisting}
<NumNodes>
\end{lstlisting}

followed by \texttt{NumNodes} lines:

\begin{lstlisting}
x    y    T    Q
\end{lstlisting}

\begin{itemize}[nosep]
  \item \texttt{x}, \texttt{y}: node coordinates in model units
  \item \texttt{T}: computed temperature (K)
  \item \texttt{Q}: boundary/conductor marker (int):
    \begin{itemize}[nosep]
      \item $-2$ = interior node (no boundary condition)
      \item $-1$ = fixed temperature boundary
      \item $0 \ldots N-1$ = conductor index
    \end{itemize}
\end{itemize}

Node indices are 0-based throughout the solution section.

% ----------------------------------------------------------------------
\subsection{Mesh Elements}

\begin{lstlisting}
<NumElements>
\end{lstlisting}

followed by \texttt{NumElements} lines:

\begin{lstlisting}
p0    p1    p2    labelIdx
\end{lstlisting}

\begin{itemize}[nosep]
  \item \texttt{p0}, \texttt{p1}, \texttt{p2}: vertex node indices (0-based)
  \item \texttt{labelIdx}: block label index (0-based, into the block
        label list from the \texttt{.feh} header)
\end{itemize}

% ----------------------------------------------------------------------
\subsection{Conductor Solution Data}

\begin{lstlisting}
<NumConductors>
\end{lstlisting}

followed by \texttt{NumConductors} lines:

\begin{lstlisting}
T    q
\end{lstlisting}

\begin{itemize}[nosep]
  \item \texttt{T}: temperature on the conductor (K).  For
        prescribed-temperature conductors, this is the prescribed
        value; for prescribed-flow conductors, this is the computed
        value.
  \item \texttt{q}: total heat flow on the conductor (W).  For
        prescribed-flow conductors, this is the prescribed value; for
        prescribed-temperature conductors, this is the computed value.
\end{itemize}

% ======================================================================
\section{Index Conventions}

\begin{itemize}
  \item \textbf{Property indices} (materials, boundaries, conductors,
        point properties) are \textbf{1-based} in the file.  An index
        of 0 means ``no property assigned.''  The parser converts to
        0-based internally.

  \item \textbf{Node indices} (in segments, arcs, and the solution
        section) are \textbf{0-based}.

  \item \textbf{Block label indices} in the solution element list are
        \textbf{0-based}, referencing the block label array from the
        \texttt{.feh} header.
\end{itemize}

% ======================================================================
\section{Unit Conversion}

The \verb|[LengthUnits]| tag determines the coordinate system.  All
coordinates in the file are in the specified units.  Conversion to
meters:

\begin{center}
\begin{tabular}{ll}
\toprule
Unit & Meters per unit \\
\midrule
inches      & 0.0254 \\
millimeters & 0.001 \\
centimeters & 0.01 \\
meters      & 1.0 \\
mils        & $2.54 \times 10^{-5}$ \\
microns     & $10^{-6}$ \\
\bottomrule
\end{tabular}
\end{center}

\end{document}
